\clearpage
\SourceLeaf{143}
\begin{center}
[III.]

{[ON THE HISTORY OF THE DISCOVERY OF THE METHOD OF\\
LEAST SQUARES.]}
\end{center}

\begin{center}
Allgemeine Geographische Ephemeriden. Edited by F.\ \textsc{von Zach}.\\
Fourth volume. 4th issue. October 1799. P.\,378. Weimar, 1799.
\end{center}

Allow me to report a misprint noticed in the July issue\ of the A.\,G.\,E. On p.\,xxxv of the introduction, in the value given for the arc between the Panthéon and Évaux, for 76545,74 read 76145,74. The sum is correct, and the error can lie only in this particular segment\HistoricalNote{*)}{*)}{The misprint is correctly identified, and can also be seen from the accompanying decimal degree $2^{0}\!,66868$. v.\,Z.}. I discovered this error by applying my method, of which I have given you a sample\HistoricalNote{**)}{**)}{More on this another time. v.\,Z.}, to determine the ellipse from these four measured segments alone, and found a flattening of $\frac{1}{150}$. After correcting that error, I found $\frac{1}{187}$, and the entire quadrant to be 2565006 modules (that is, disregarding the degree in Peru). The difference between $\frac{1}{150}$ and $\frac{1}{187}$ is not particularly significant in this case, since the endpoints are too close together.

\bigskip
\noindent Braunschweig, 24 Aug.\ 1799. \hfill C.\,F.\,\textsc{Gauss}.

